% Pursuing Stacks (cotes 134-2 to 134-8): archive cote 134-4, site page 14 (PDF page 15), typescript page 102.
% typescript. Draft compared: ps3.tex:691-761.
% AI-transcribed by Claude Opus 5.5 (claude-opus-5-5, adaptive thinking, effort high), Message Batches API, batch ps1 (msgbatch_01Rv5hCaNTtUt4JHx1yqoHPL), 2026-11-12; not yet checked by eye.
\scanpage{134-4/14}
\mspage{102}
indices.  T\over{\illegible{1}}{h}e dual statement so to say,  when X is decomposed as a sum
$X = \coprod X_i$ ,  is valid too,  provided taking products with I is distributive
with respect to \inserted{finite} direct sums.  Thus we get that $M \to M_{\underline{I}}$  commutes with \deleted{ar-}%
\deleted{bitrary} \inserted{finite} \xout{xxxx} products,  and with \deleted{arbitrary} \inserted{finite} sums too provided they are dis-%
tributive with respect to multiplication with any object  (or \inserted{with} I \hand{only},
which would be enough).

The notion of $\underline{I}$-homotopy in the wide sense between $f,g\colon X \rightrightarrows Y$ can
be interpreted in terms of strict $\underline{I}'$-homotopy with variable $\underline{I}'$ ,  as fol-%
lows,  provided we make some mild extra assumptions on (M,W), namely:

a)  (Just for memory)  M is stable under finite products.

b)  M is stable under amalgamated sums  $I \amalg_e J$  under the final ob-%
ject e (I and J endowed with sections over e).

c)  If moreover I and J are aspheric over e ,  then so is $I \amalg_e J$.

Conditions \over{\unsure{a}}{b}) and \over{\unsure{b}}{c}) give a means of constructing new homotopy inter-%
vals\inserted{$\underline{K}$},  by amalgamating two \inserted{homotopy} intervals  $\underline{I}$ and $\underline{J}$ ,  using as sections \xout{fxx}
of $\underline{I}$ and $\underline{J}$ for making the amalgamation,  either $\hand{$\delta$}_0$ of $\hand{$\delta$}_1$ ,  which gives
four ways of amalgamating  -  of course we take as endpoints of the amal-%
gamated\xout{x} interval, \inserted{the} sections over e coming from the two endpoints of
$\underline{I}$  (giving rise to $\hand{$\delta$}_0$\inserted{for K}) and $\underline{J}$  (giving rise to $\hand{$\delta$}_1$ for $\underline{K}$) which have not
been ``used up'' in the amalgamation.  Maybe the handiest convention is to
define \uline{the}  amalgamated interval $\underline{I}\vee\underline{J}$ ,  without any ambiguity of choice,
as being
\[ \xout{$\underline{K}$}\ \underline{I}\vee\underline{J} = (I,\hand{$\delta$}_1^{\hand{$I$}})\amalg_e(J,\hand{$\delta$}_0^{\hand{$\unsure{J}$}}) \qquad \text{endowed with  the two sections coming from } \hand{$\delta$}_0^{\hand{$I$}}\colon e\to I \ , \ \hand{$\delta$}_1^{\hand{$J$}}\colon e\to J \ , \]
and defining the three other choices \xout{bxx} in terms of this operation,  by
replacing one or two among the summands $\underline{I}$, $\underline{J}$  by the ``opposite'' interval
$\overset{\hand{$\vee$}}{\underline{I}}$  or  $\overset{\hand{$\vee$}}{\underline{J}}$ .  T\over{\illegible{1}}{h}e operation of amalgamation of intervals,  and likewise of
homotopy intervals,  just defined,  is clearly associative up to canonical
isomorphism,  and we have a canonical isomorphism of intervals
\[ (\underline{I}\vee\underline{J})^{\hand{$\vee$}} \simeq \overset{\hand{$\vee$}}{\underline{I}}\vee\overset{\hand{$\vee$}}{\underline{J}} \ . \]
I forgot to check,  for amalgamation of homotopy intervals,  the condition
b) of page 59 namely $e_0 \wedge e_1 = \emptyset_M$  (which has not so far played any role,
anyhow).  To get this condition,  we'll have to be slightly more specific
in condition b) above on M of existence of the relevant amalgamations,
by demanding (as suggested of course by the visual intuition of the situa-%
tion)  that I and \over{\illegible{1}}{J}  should become subobjects of the amalgamation K,  and
their intersection should be reduce\over{\illegible{1}}{d} to \over{\illegible{1}}{t}he tautological part e of it.
M\over{\illegible{1}}{o}re relevant still for the use we have in mind is to demand that those
amalgamations should commute to products b\over{\illegible{1}}{y} an arbitrary element X of M.

\margin{14}
